\section{Introducción}
\subsection{Descripción del problema}

En los últimos años, la inteligencia artificial ha avanzado hasta resolver tareas que hace una década parecían prohibitivas. La detección de objetos en tiempo real, el análisis de imágenes médicas y la traducción automática son problemas que hoy tienen soluciones competitivas, a veces superiores al desempeño humano en condiciones controladas. Sin embargo, no todo el territorio está cubierto: la música resiste, y lo hace por razones que no se resuelven añadiendo más datos o más capas, como la percepción, el contexto y las decisiones interpretativas, que no caben fácilmente en una función de pérdida.

La transcripción musical automática (AMT) sigue siendo un reto, ya que va más allá del procesamiento y análisis de audio para su posterior transcripción. En el caso de los instrumentos polifónicos, se debe tener en cuenta la superposición espectral que se produce al formar acordes, lo que conlleva consideraciones más detalladas, como el material de las cuerdas. En la guitarra acústica con cuerdas de nailon, por ejemplo, se presenta una menor inarmonicidad y un \textit{decay} más pronunciado que en las cuerdas metálicas, lo que altera la distribución de armónicos en el espectrograma.

Además de las condiciones físicas del instrumento, se deben tener en cuenta las distintas técnicas de ejecución. Por ejemplo, un Am7 puede sonar perceptiblemente distinto según se toque con dedos o púa, en distintas posiciones del mástil, con vibrato, \textit{muting} parcial o cambios de dinámica. De igual forma, conceptos de la teoría musical, como la enarmonía, introducen otra capa de complejidad. Do\# y Re$\flat$ comparten frecuencia fundamental, pero representan acordes distintos según el contexto armónico, una distinción que el espectrograma por sí solo no puede resolver y que requiere contexto adicional o intervención humana para mejorar su precisión.
